\documentclass[11pt]{article}
\usepackage[margin=1in]{geometry}
\usepackage{amsmath,amssymb,amsthm,hyperref}
\newtheorem{lemma}{Lemma}
\newtheorem{theorem}[lemma]{Theorem}
\newtheorem{proposition}[lemma]{Proposition}
\begin{document}
\setlength{\emergencystretch}{2em}
\title{A two-adic residue reduction for ordinary support five/seven}
\author{}\date{}\maketitle
We prove ordinary consecutive coprimality outside two lower-index
classes modulo sixteen. At every index, a possible nonconstant
common gcd is an irreducible imaginary quadratic.
The full support-five/seven SFC family remains open.
Let
\[
 P_n(t)=\sum_{k=0}^n\binom nk(5n+1)_{2k}t^k,\quad
 \gamma=5n+1,\quad M=7n+1,\quad s=2(n+1).
\]
Every valuation is normalized by $v_2(2)=1$, in a complete
algebraically closed extension of $\mathbb Q_2$.
Set $B=B_{n,\gamma,2}$, $C=B_{n+1,\gamma,2}$, always at the same
$\gamma$, with
$B_0=1$ and $B_{j+1}=(z+2j+\gamma+1)B_j-zB_j'$.
Write $J=J_0(z^2)+zJ_1(z^2)$ for its two cyclic components.

\begin{lemma}[The exact shift-five condition]\label{lem:shift}
Put $A_5=(M+2)_5$ and
\[
 U=A_5+3s(M+4)^2u+su^2,\qquad
 W=su[(M+4)(M+5)(M+6)+(M+2n+8)u].
\]
Every common root $t$ of $P_n,P_{n+1}$ gives
$u=-1/t\ne0$ with
\[
 U(u)C_i(u)+W(u)B_i(u)=0\quad(i=0,1).
\]
\end{lemma}
\begin{proof}
The full rising-jet concomitant in the proof of \texttt{thm:mixed}
of the \href{LOWER_TWO_PRIME_GAP_SFC.tex}{general mixed-rank note}, at $d=2$ and $R=0$,
is $B(\pm z)(E)_2R-C(\pm z)ER=0$, where
$R=z^{2n}P_n(-z^{-2})$ and $E=M-zD_z$.
For the upper transformed moment,
\[
 (\gamma)_5z^{2(n+1)}T_{n+1,\gamma+5,2}(-z^{-2})
              =[u(E)_5-(E)_7]R,\qquad u=z^2.
\]
This follows from Pascal's identity and
$(E)_h$ acting on the index-$k$ term by $(\gamma+2k)_h$.
The upper moment operator is therefore $u(E)_5-(E)_7$.
Modulo the lower cubic equation it is $qR-WER-U(E)_2R$
for a polynomial $q$. This identity is verified completely by
the following formal recurrence, with $\gamma=M-2n$.
For a jet coefficient vector $(r,e,k)$ in basis $R,ER,(E)_2R$,
application of $E$ gives
\[
 (-2ur'-\gamma uk,\quad r-e-2ue'+uk,\quad e+Mk-2uk'),
\]
where primes mean $d/du$. Starting at $(1,0,0)$, apply $E+j$
successively to obtain $(E)_hR$ for $h=5,7$.
The last two coefficients of $u(E)_5-(E)_7$ are exactly $-W,-U$.

At an ordinary point $z\ne0$, the actual nonzero cubic solution jet
cannot have $R=ER=(E)_2R=0$. Thus at a common root the nonzero
pair $(ER,(E)_2R)$ solves both the upper relation and each
concomitant relation. Their two by two determinant is
$UC(\pm z)+WB(\pm z)$, which must vanish.
Characteristic-zero Fourier inversion gives the two cyclic equations.
No division by $U$, $B$ or $R'$ is used.
\end{proof}

Every lower root has $L=v_2(u)\ge1$, because its coefficient at
degree $k$ is divisible by $2^k$.
At such $u$, the reductions of $B,C$ are both $z$ if $n$ is even,
and both $1+z$ if $n$ is odd, on $k[z]/(z^2)$.
Indeed the fixed reduced operator is
$z-zD_z+\bar\gamma+1$ and its positive powers applied to $1$
are respectively $z$ and $1+z$. In particular these columns never
reduce to zero.

\begin{theorem}[Fourteen of sixteen classes]\label{thm:residues}
For every $n\ge1$ with $n\not\equiv11,13\pmod{16}$,
$\gcd(P_n,P_{n+1})=1$ over $\mathbb Q$ and $\mathbb C$.
Equivalently, this argument leaves only terminal orders
$N\equiv12,14\pmod{16}$ for ordinary support $\{5,7\}$.
\end{theorem}
\begin{proof}
First let $n$ be even and put $m=v_2(A_5)\ge3$.
The constant, linear and quadratic coefficients of $U$ have
valuations $(m,1,1)$.
The linear coefficient of $W$ has valuation at least two and
its quadratic coefficient has valuation one, since
$M+2n+8=9(n+1)$ is odd.
If $L<m-1$, the linear term of $U$ is uniquely least, and
$W/U$ has positive valuation. The same holds with the constant
term least if $L>m-1$. Dividing the cyclic equation by $U$
would give $\bar C=0$. Hence $L=m-1$.
With $q_0=2^{m-1}$,
\[
 F_n(X)=2^{-m}[U(q_0X)C_1(q_0X)+W(q_0X)B_1(q_0X)]
\]
is an integral finite polynomial with reduction $1+X$.
Its difference quotient at any two integral arguments is a unit,
so it has at most one integral zero. Every conjugate of an actual
common root gives a unit zero via $X=-1/(q_0t)$.
A nonconstant rational common factor must therefore be linear,
forcing a real common root. The even-index moment polynomial
is positive on the real axis, a contradiction.

Now let $n$ be odd, put $N=n+1$ and $\ell=v_2(N)$.
For $\ell=1$, set $b=v_2(M+4)=v_2(7n+5)\ge2$.
The valuations of the three $U$ coefficients are
$(b+2,\,2+2b,\,2)$; the two $W$ coefficients have valuations
$(b+3,\,3)$. Comparing terms exactly as above forces $L=b/2$.

Whenever $n$ is odd, write $\gamma=2g$. The exact coefficient valuation is
\[
 v_2\!\left(\binom nk(\gamma)_{2k}\right)
   =k+v_2\binom nk+v_2(g)_k.                       \tag{1}
\]
For $n\equiv1\pmod4$, $g$ is odd and the coefficient at
$k=1$ has valuation one. For all $k\ge2$
the right side of (1) is STRICTLY greater than
$1+3(k-1)/2$: for even $k$, use $v_2(g)_k\ge k/2$;
for odd $k\ge5$, two even factors include a multiple of four,
giving $v_2(g)_k\ge(k-1)/2+1$; for $k=3$ use
$v_2\binom n3=v_2(n-1)-1\ge1$.
Thus at $L=3/2$ the linear term of $P_n(t)$ is uniquely least.
This excludes $b=3$, namely $n\equiv5\pmod{16}$.
For $b=2$, namely $n\equiv1\pmod8$, one has $L=1$.
The actual lower polynomial $P_n(X/2)$ is integral with reduction
$1+X$, since $v_2(g)_k\ge\lfloor k/2\rfloor\ge1$ for $k\ge2$.
All common-root conjugates would give its unique integral zero,
again forcing a real common root and a contradiction.

For $\ell\ge3$, the valuations of the $U$ coefficients are
$(\ell+3,\ell+3,\ell+1)$, and those of $W$ are
$(2\ell+2,2\ell+1)$. The only possible balance is $L=1$.
But $\gamma=5N-4$ is divisible by four. With $g=\gamma/2$
even, (1) gives valuations at least $k+\lceil k/2\rceil>k$
for all $k\ge1$, so at $L=1$ the constant term is uniquely least.
This excludes $n\equiv7\pmod8$.

Finally let $\ell=2$ and $a=v_2(M+2)=v_2(7n+3)\ge3$.
When $a=3$, namely $n\equiv3\pmod{16}$, the $U$ valuations
are $(6,5,3)$ and the $W$ valuations are $(6,5)$.
Their only possible balance is $L=3/2$.
Here $\gamma=5N-4$ is divisible by eight. Thus $g=\gamma/2$
is divisible by four, and
\[
 v_2(g)_k\ge\lceil k/2\rceil+\lceil k/4\rceil>k/2.
\]
All nonconstant terms of $P_n(t)$ then have positive valuation
at $L=3/2$, contradicting its constant term one.

The unexcluded odd cases are exactly $n=11,13\pmod{16}$.
For the real-root contradiction used above, even $j$ gives
\[
 P_j(t)=\frac1{\Gamma(5j+1)}
 \int_0^\infty x^{5j}(1+tx^2)^j e^{-x}\,dx>0
 \quad(t\in\mathbb R).
\]
These are uniform valuation arguments, not an initial-index search.
\end{proof}
\begin{proposition}[Two local candidates and a quadratic residual obstruction]
\label{prop:quadratic-obstruction}
For every $n\equiv11\pmod{32}$, every common root gives $v_2(u)=2$.
The two polynomials
\[
 F_i(X)=2^{-7}\bigl[U(4X)C_i(4X)+W(4X)B_i(4X)\bigr],
 \qquad i=0,1,
\]
belong to $\mathbb Z[X]$ and both reduce to $1+X+X^2$.
There are at most two distinct common roots over $\mathbb C_2$.
The rational gcd of $P_n,P_{n+1}$ is consequently either constant
or an irreducible quadratic.

Every integral polynomial linear combination of $F_0,F_1$ has
reduction divisible by $1+X+X^2$. No such combination has a nonzero
affine reduction without a further division or cancellation step.
These conclusions do not assert that a common root exists.
\end{proposition}
\begin{proof}
Here $\ell=2$, $v_2(M+2)=4$, and $v_2(A_5)=7$.
The three coefficients of $U$ have valuations $(7,5,3)$,
and the two coefficients of $W$ have valuations $(6,5)$.
For $1\le L<2$ the quadratic term of $U$ is uniquely least and
both terms of $W$ have strictly larger valuation.
For $L>2$ its constant term is uniquely least and again $W/U$
has positive valuation. The nonzero reduced column of $C$
excludes both cases. Thus a common root has $L=2$.

At $u=4X$, all three terms of $U(4X)$ have coefficient valuation
seven; both terms of $W(4X)$ have strictly larger coefficient
valuation. Both components of $C(4X)$ reduce to one.
This proves the stated integrality and whole polynomial reduction.

The polynomial $Q=1+X+X^2$ has two distinct roots in the algebraic
closure of $\mathbb F_2$, and $Q'=1$ there. In either residue disk,
the divided difference of $F_0$ at any two distinct integral points
reduces to $Q'(\rho)=1$. It is a unit, so the disk contains at most
one zero. Any integral zero must reduce to a root of $Q$.
Thus $F_0$ has at most two integral zeros.
The map $t\mapsto-1/(4t)$ is injective on nonzero roots, so there are
at most two distinct actual common roots.

For the multiplicity assertion we import Theorem
\texttt{thm:simple} of the
\href{CUBIC_ADJOINT_STABILITY_SFC.tex}{quadratic-gap simplicity note}.
It applies at $\gamma=5n+1\ge n-2$, hence $P_n$ is squarefree.
Therefore the rational gcd has degree at most two.
It has no real root by the even-moment positivity proved above,
so if nonconstant it must be an irreducible quadratic.

Finally every integral polynomial combination reduces to a multiple
of $Q$. This multiple is either zero or has degree at least two.
The residual quadratic is irreducible over $\mathbb F_2$.
The obstruction concerns these normalized rows before any further
division; it does not preclude a higher-precision coupled argument.
No common ordinary root or SFC counterexample has been produced.
\end{proof}


\begin{theorem}[A global quadratic common-factor bound]
\label{thm:global-quadratic-bound}
For every $n\ge1$, the rational gcd of $P_n,P_{n+1}$ is either
constant or an irreducible quadratic of negative discriminant.
For every odd $n$, put $\ell=v_2(n+1)$ and
\[
 J_n(u)=2^{-\ell-1}
       [U(u)C_1(u)+W(u)B_1(u)].
\]
This is an integer polynomial whose whole reduction satisfies
\[
 \bar J_n(u)=u^2\bar C_1(u),\qquad \bar C_1(0)=1.
\]
Thus there are at most two distinct actual common roots over
$\mathbb C_2$ at every odd lower index, without restricting the
valuation of $u$ to one annulus.
\end{theorem}
\begin{proof}
Let $n$ be odd and set $N=n+1$, so $M=7N-6$ is even and $\ell\ge1$.
Write $a=v_2(M+2)\ge1$ and $b=v_2(M+4)\ge1$.
The other even factor of $A_5$ is $M+6=7N$, of valuation $\ell$;
the remaining two factors are odd.
Consequently the three coefficients of $U$, in ascending order,
have valuations
\[
 \ell+a+b,\qquad \ell+1+2b,\qquad\ell+1.
\]
The two nonzero coefficients of $W$ have valuations
\[
 2\ell+1+b,\qquad2\ell+1,
\]
using $M+2n+8=9N$.
It follows that the constant and linear coefficients of
$2^{-\ell-1}U$ are even, its quadratic coefficient is odd, and
every coefficient of $2^{-\ell-1}W$ is even.
Since $B_1,C_1$ have integer coefficients, this proves integrality
of $J_n$ and the whole residual identity
$\bar J_n=u^2\bar C_1$.
At $u=0$, the fixed reduced auxiliary operator is
$A=z-zD_z+1$ on $k[z]/(z^2)$.
It sends $1$ to $1+z$ and $z$ to zero, so every positive power
on $1$ is $1+z$. Hence $\bar C_1(0)=1$, and the order of
$\bar J_n$ at zero is exactly two.

Here is an elementary root count in this residue disk.
If an integral polynomial has a root $\alpha$ in the maximal ideal,
division by the monic factor $u-\alpha$ preserves integrality:
starting with its integral leading coefficient, synthetic division
recursively adds an integral coefficient and $\alpha$ times the
previous quotient coefficient.
After factoring off $r$ such roots, counted with multiplicity,
reduction is therefore divisible by $u^r$.
The remaining quotient has nonzero reduction, since the original
reduction was nonzero and the reduced product is $u^r$ times that
quotient. A polynomial with residual order exactly two thus has
at most two roots in the maximal ideal, counted with multiplicity.

Every actual common root maps injectively into that disk by
$u=-1/t$, as proved before Theorem~\ref{thm:residues}, and is a root
of $J_n$ by Lemma~\ref{lem:shift}.
There are at most two distinct common roots.
For the rational gcd degree, import \texttt{thm:simple} of the
\href{CUBIC_ADJOINT_STABILITY_SFC.tex}{quadratic-gap simplicity note}
at $\gamma=5n+1>0$, $\gamma\ge n-2$.
It makes $P_n$ squarefree, so the gcd has degree at most two.
The even-index member has no real root by the positive integral
in Theorem~\ref{thm:residues}.
A nonconstant gcd is therefore an irreducible quadratic with
negative discriminant.
For even $n$, Theorem~\ref{thm:residues} already proves coprimality.
\end{proof}

\paragraph{Evidence and limits.}
The checker verifies the symbolic shift-five jet reduction, literal
valuation identities and complete residue partition, full even-index
scalar reductions, and quadratic residual profiles.
It also checks the full normalized odd-index component identity and
residual order two for the global degree bound.
The written valuation and synthetic-division arguments supply the
unbounded quantifiers; finite controls do not extrapolate them.
The support $\{5,7\}$ remains open on the two stated residue rays,
where a possible common factor is now uniformly confined to an
irreducible imaginary quadratic.
Internal replay does not establish independent mathematical review,
formal verification or priority.
\end{document}
